% ============================================================
% MÓDULO: SISTEMAS DE GESTIÓN EMPRESARIAL (ERP)
% ============================================================
\clearpage
\section*{MÓDULO: SISTEMAS DE GESTIÓN EMPRESARIAL (ERP)}
\subsection*{Odoo, SAP y Salesforce}

\begin{recuadro}
\textbf{Objetivo del módulo:} comprender qué es un ERP, para qué lo usan las empresas y cómo
se registran en él las operaciones de compra, stock, producción, mantenimiento y venta,
trabajando con un sistema real (Odoo).
\end{recuadro}

% ------------------------------------------------------------
\section*{1. ¿QUÉ ES UN ERP?}

Un \textbf{ERP} (\textit{Enterprise Resource Planning}, planificación de recursos
empresariales) es un sistema informático que \textbf{integra en una sola base de datos} toda la
información de una empresa: compras, inventario, producción, ventas, facturación, contabilidad,
personal y mantenimiento.

Sin un ERP, cada sector suele llevar su propia planilla: el pañol anota sus salidas en un
cuaderno, compras usa un Excel, administración factura en otro sistema. La información se
duplica, se contradice y se pierde. Con un ERP, cuando el pañol entrega un componente para una
orden de trabajo, el stock se descuenta automáticamente, el costo se imputa a esa orden y, si
se llega al stock mínimo, se genera la solicitud de compra.

\subsubsection*{1.1. Ventajas}
\begin{itemize}
  \item Una \textbf{única fuente de datos}: todos trabajan sobre la misma información.
  \item \textbf{Trazabilidad:} se sabe quién hizo qué y cuándo (fundamental en aviación).
  \item \textbf{Automatización} de tareas repetitivas (reposición, facturación, alertas).
  \item \textbf{Informes} en tiempo real para tomar decisiones.
\end{itemize}

\subsubsection*{1.2. Módulos típicos}
\begin{center}
\small
\begin{tabular}{|p{0.25\linewidth}|p{0.65\linewidth}|}
\hline
\textbf{Módulo} & \textbf{Función} \\ \hline
Inventario / Almacén & Productos, ubicaciones, entradas, salidas, lotes, números de serie. \\ \hline
Compras & Solicitudes de cotización, órdenes de compra, recepción de mercadería. \\ \hline
Ventas & Presupuestos, pedidos de clientes, entregas. \\ \hline
Fabricación (MRP) & Listas de materiales, órdenes de producción o de trabajo. \\ \hline
Mantenimiento & Equipos, mantenimiento preventivo y correctivo. \\ \hline
Contabilidad y facturación & Facturas, cobros, pagos, impuestos. \\ \hline
CRM & Seguimiento de clientes y oportunidades de venta. \\ \hline
Recursos Humanos & Empleados, asistencia, vacaciones. \\ \hline
\end{tabular}
\end{center}

% ------------------------------------------------------------
\section*{2. LOS PRINCIPALES SISTEMAS}

\subsubsection*{2.1. SAP}
Empresa alemana, líder mundial en ERP para grandes compañías. Su producto actual es
\textbf{SAP S/4HANA}. Es muy completo y robusto, pero su implementación es larga y costosa.
Lo usan aerolíneas, fabricantes y grandes industrias; conocerlo es un valor agregado en el CV.

\subsubsection*{2.2. Salesforce}
Plataforma en la nube especializada en \textbf{CRM} (gestión de la relación con el cliente):
ventas, atención al cliente y marketing. No es un ERP completo, pero suele integrarse con uno.

\subsubsection*{2.3. Odoo}
ERP de origen belga, \textbf{de código abierto} y modular: se instalan solo las aplicaciones
que se necesitan. Tiene dos ediciones:
\begin{itemize}
  \item \textbf{Community:} gratuita y de código abierto (licencia LGPLv3). Incluye, entre
        otras, Inventario, Compras, Ventas, Fabricación, Mantenimiento, Reparaciones,
        Facturación, Sitio web y Contactos. Se instala en un servidor propio.
  \item \textbf{Enterprise:} paga. Agrega aplicaciones avanzadas (por ejemplo, Calidad, PLM,
        Studio, contabilidad completa y la localización argentina oficial con
        facturación electrónica ARCA), la aplicación móvil, hosting y soporte.
\end{itemize}
Para el uso educativo y para pymes, Odoo Community es ideal: es gratis, se puede instalar en
cualquier PC o servidor y además existen miles de módulos libres de la comunidad (OCA,
\textit{Odoo Community Association}), entre ellos módulos de mantenimiento y de localización
argentina.

\begin{center}
\small
\begin{tabular}{|p{0.18\linewidth}|p{0.22\linewidth}|p{0.22\linewidth}|p{0.22\linewidth}|}
\hline
 & \textbf{SAP} & \textbf{Salesforce} & \textbf{Odoo} \\ \hline
Tipo & ERP integral & CRM en la nube & ERP modular \\ \hline
Destinatarios & Grandes empresas & Áreas comerciales & Pymes y medianas \\ \hline
Licencia & Paga & Paga (suscripción) & Community libre / Enterprise paga \\ \hline
Instalación & Nube o servidor propio & Nube & Nube o servidor propio \\ \hline
\end{tabular}
\end{center}

% ------------------------------------------------------------
\section*{3. CONCEPTOS CLAVE DE INVENTARIO EN UN ERP}

\begin{itemize}
  \item \textbf{Producto:} cada artículo, con su código interno (SKU), unidad de medida, costo
        y proveedor. Puede ser \textit{almacenable} (se controla stock), \textit{consumible}
        (no se controla) o \textit{servicio}.
  \item \textbf{Almacén y ubicaciones:} el almacén se divide en ubicaciones (estantería, cajón,
        gaveta). Ejemplo: \texttt{PAÑOL/Est-A/Caj-03}.
  \item \textbf{Movimiento de stock:} toda entrada, salida o transferencia queda registrada
        con fecha, usuario y documento de origen.
  \item \textbf{Lotes y números de serie:} permiten la trazabilidad de cada unidad. En
        aviación, cada componente instalado debe poder rastrearse hasta su certificado de
        origen.
  \item \textbf{Regla de reabastecimiento (mín./máx.):} cuando el stock baja del mínimo, el
        sistema propone comprar hasta el máximo (el ``Kanban'' digital).
  \item \textbf{Inventario físico:} conteo real para ajustar las diferencias con el sistema.
  \item \textbf{Lista de materiales (BoM):} los componentes necesarios para fabricar o reparar
        algo.
\end{itemize}

% ------------------------------------------------------------
\section*{4. ACTIVIDADES DEL MÓDULO}

\begin{enumerate}
  \item \textbf{Del cuaderno al ERP:} a partir de una situación del taller (un pedido de
        componentes para reparar un equipo), dibujar el recorrido de la información en papel y
        luego en un ERP. ¿Qué pasos desaparecen? ¿Qué errores se evitan?
  \item \textbf{Exploración de Odoo:} usando una base de demostración, crear un producto, un
        proveedor, una orden de compra, recibir la mercadería y verificar el stock.
  \item \textbf{Investigación:} buscar avisos laborales que pidan SAP, Odoo o Salesforce. ¿Qué
        puestos son? ¿Qué conocimientos piden?
\end{enumerate}
